\chapter{遥感与航拍}

\paragraph*{本章阅读锚点}
遥感与航拍里的 3DGS 不能只按新视角合成质量来读。原始 3DGS 的锚点是由 SfM 或 COLMAP 给出相机和稀疏点，再用可微 splatting 优化一组带颜色、不透明度、尺度和协方差的高斯；这一套在近景照片上很有效，但到了测绘、作物表型和卫星重建，任务目标往往变成可量测的正射图、三维坐标、植株器官实例或 DSM 高程。于是本章的方法大多不是简单把 3DGS 换成另一个渲染器，而是在 3DGS 外面接入摄影测量约束、遥感相机模型、语义分割、正射投影或人工量测流程，让高斯场从“看起来像”进一步服务于“可定位、可分割、可计量”。

\section{文化遗产与测绘}

文化遗产和测绘任务原来的锚点是摄影测量：航片先做空三和姿态估计，再由 MVS 或 DSM 生成网格、点云、DOM 或 TDOM。3DGS 带来的变化是把中间表示换成高保真、可实时浏览的高斯场，但这也带来一个新问题：高斯不是显式网格，不能直接按传统方式查询表面坐标或做遮挡检测。因此这一节的两篇工作分别从两个方向补上测绘语义：A-TDOM 把正射图生成嵌入在线 3DGS 训练过程，Point Measurement 则把 3DGS 当作交互式视觉参考，再回到经典空间前方交会求三维点。



% METHOD: 2509.12759
\paragraph{A-TDOM（在线 TDOM 生成）（2025.09）}
A-TDOM~\cite{2509_12759} 继承的是“航测空三加 DSM 生成 TDOM”的测绘路线，同时把表示层换成原始 3DGS。传统 TDOM 的瓶颈在于需要先完成影像定向、DSM 构建、正射纠正和遮挡检测；离线 3DGS 虽然能重建场景，却通常等所有图像到齐后再训练，无法在无人机采集过程中逐步看到正射成果。A-TDOM 的改动是把 On-the-Fly SfM、渐进式 3DGS 训练和正射 splatting 串成一个在线循环，让每张新航片进入系统后都更新相机、扩展高斯场并刷新 TDOM。

实际流程是先等待初始的若干张图像完成在线 SfM，得到稳定的相机位姿和稀疏点云，再按原始 3DGS 初始化一个较小的高斯场。之后每来一张新图像，系统先注册该图像并微调已有位姿，再用重投影的 Delaunay 三角剖分确定当前图像中与已有模型有足够重叠的关键区域，避免在低重叠或未可靠定位区域盲目扩张。对确实新出现或重建粗糙的区域，它执行 Gaussian Sampling and Integration：在这些区域内采样新的高斯，把它们并入当前高斯场，然后只给新区域和受影响区域分配更集中的优化迭代。训练损失仍围绕 3DGS 的图像重建项和结构相似度项展开，但迭代预算、学习率和可见性更新被改成适合流式输入的调度。最后，系统从更新后的高斯场做正射渲染，项目文件中的在线训练脚本也体现了这一点：工程入口围绕 On-the-Fly SfM 并发、渐进训练、TDOM 相机和中间 TDOM 渲染结果组织，而不是单次离线训练。它的价值在于把 3DGS 从“航片处理后的可视化模型”推进成“采集时就能更新的测绘产品生成器”。
\textbf{方法信息：} 代码：\url{https://github.com/xywjohn/a-tdom}。



% METHOD: 2603.24716
\paragraph{3DGS 点量测（渲染辅助的空间前方交会）（2026.03）}
Accurate Point Measurement in 3DGS~\cite{2603_24716} 的锚点不是某个新的高斯训练算法，而是传统摄影测量中的空间前方交会。它观察到，工程量测常常卡在网格或稠密点云的几何误差上：MVS 网格会平滑细杆、锐角和反光表面，直接在网格上点选坐标时，坐标精度受显式几何质量限制。该方法的改变是把 3DGS 的角色降到“高保真、多视角、连续浏览的量测界面”，不从高斯本身直接读表面点，而是让用户在若干渲染视角中点选同一个语义点，再用相机中心和像素射线做多视图最小二乘交会。

因此它的流程很清楚：输入是已有相机位姿的 UAV 3DGS 场景，用户在交互式查看器里从多个视角点击同一个 congruent point；系统把每次点击反投影成一条三维视线，多条射线不要求严格相交，而是通过最小二乘求离所有射线最近的三维点，并可用不确定性椭球表示量测稳定性。公开实现是一个 Web/Cesium 工具，README 中的使用流程对应“加载 3DGS 模型、点击多视图射线、生成点线面、导出量测结果”；它还把 3DGS 转成可流式浏览的 3D Tiles/GLB 资源，使工程用户不必进入传统立体测图工作站。与 A-TDOM 相比，这篇工作不改变 3DGS 的训练，而是改变量测范式：它承认高斯场不是严谨网格，却利用 3DGS 保留细节和自由浏览的优势，把最终坐标计算交给成熟的摄影测量几何。实验中厘米级 RMSE 和细结构优势说明，这类“在渲染上识别、在几何上求交”的路线适合测绘中需要人工确认目标点的场景。
\textbf{实现细节：} TOrtho-Gaussian（2024.11） 首次提出了通过正交投影 3DGS 高斯核生成 TDOM 的方法。与传统摄影测量管线需要先生成数字表面模型（DSM）再进行遮挡检测和正射纠正不同，该方法直接将 3DGS 的 3D 高斯基元通过正交投影映射到地面平面，利用高斯的不透明度加权自然处理遮挡关系——距离相机更近的高斯具有更高的渲染权重，自动遮挡后方高斯。采用分而治之策略处理大规模区域。与锚点 TOrtho-Gaussian 将 3DGS 高斯基元直接用于正交投影渲染 TDOM 的产品输出不同，Deng and Qin（2026.03） 将 3DGS 的角色从渲染引擎转变为高保真视觉参考——利用其无缝新视角合成能力作为操作者瞄准测量射线的交互界面，而非生成任何正射影像产品。这一设计动机在于：传统立体测量需要专业摄影测量工作站和密集匹配生成的网格模型，而网格表面的粗糙采样在薄结构和锐角处无法精确表示尖锐几何特征，导致直接在网格上的点选测量精度受限；3DGS 渲染的照级真实感使操作者能够跨视角准确识别同一语义点（"一致点"），绕过了底层网格的几何不准确性。具体算法流程为循环式多射线空间前方交会：操作者在 3DGS 场景中自由导航到第$i$个视角$V_i$，点击像素$p_i$标记一致点，渲染引擎返回当前相机位姿（相机中心$C_i = (X_i, Y_i, Z_i)^T$和从$C_i$穿过$p_i$的方向向量$d_i = (u_i, v_i, w_i)^T$），形成射线$R_i = (C_i, d_i)$；重复此过程收集$N \geq 2$条射线后，构建超定线性最小二乘系统。
\textbf{技术实现补充：} TOrtho-Gaussian（2024.11） 首次提出了通过正交投影 3DGS 高斯核生成 TDOM 的方法。与传统摄影测量管线需要先生成数字表面模型（DSM）再进行遮挡检测和正射纠正不同，该方法直接将 3DGS 的 3D 高斯基元通过正交投影映射到地面平面，利用高斯的不透明度加权自然处理遮挡关系——距离相机更近的高斯具有更高的渲染权重，自动遮挡后方高斯。采用分而治之策略处理大规模区域。与锚点 TOrtho-Gaussian 将 3DGS 高斯基元直接用于正交投影渲染 TDOM 的产品输出不同，Deng and Qin（2026.03） 将 3DGS 的角色从渲染引擎转变为高保真视觉参考——利用其无缝新视角合成能力作为操作者瞄准测量射线的交互界面，而非生成任何正射影像产品。这一设计动机在于：传统立体测量需要专业摄影测量工作站和密集匹配生成的网格模型，而网格表面的粗糙采样在薄结构和锐角处无法精确表示尖锐几何特征，导致直接在网格上的点选测量精度受限；3DGS 渲染的照级真实感使操作者能够跨视角准确识别同一语义点（"一致点"），绕过了底层网格的几何不准确性。具体算法流程为循环式多射线空间前方交会：操作者在 3DGS 场景中自由导航到第$i$个视角$V_i$，点击像素$p_i$标记一致点，渲染引擎返回当前相机位姿（相机中心$C_i = (X_i, Y_i, Z_i)^T$和从$C_i$穿过$p_i$的方向向量$d_i = (u_i, v_i, w_i)^T$），形成射线$R_i = (C_i, d_i)$；重复此过程收集$N \geq 2$条射线后，构建超定线性最小二乘系统。每条射线的几何约束来自叉积条件$(P - C_i) \times d_i = 0$，展开后得到三个方程但仅有两个线性独立，选取前两个整理为$2 \times 3$矩阵$A_i$和$2 \times 1$向量$\mathbf{b}_i$： \begin{equation*} A\_i = \begin{pmatrix} 0 \& w\_i \& -v\_i \\ -w\_i \& 0 \& u\_i \end{pmatrix}, \quad \mathbf{b}\_i = \begin{pmatrix} w\_i Y\_i - v\_i Z\_i \\ -w\_i X\_i + u\_i Z\_i \end{pmatrix} \end{equation*} $N$ 条射线堆叠为$2N \times 3$矩阵$A$和$2N \times 1$向量$\mathbf{b}$，最小二乘解为$\mathbf{x} = (A^TA)^{-1}A^T\mathbf{b}$。
\textbf{方法信息：} 代码：\url{https://github.com/gdaosu/3dgs_measurement_tool}。
\textbf{设计动机与边界：} 相对原始 3DGS 以独立高斯和图像重建误差为主的设定，Accurate Point Measurement in 3DGS -- A New Alternative to Traditional Stereoscopic-View B 把问题落在以下可实现的改动上：方法基于经典摄影测量中的空间交会原理：通过多条视线的最小二乘交点计算3D坐标。对于每条射线，定义相机中心C\_i和方向向量d\_i，构建线性方程组Ax=b求解最小二乘。3DGS的独特优势：精确渲染源视图和平滑插值中间视图，使用户可在操作3DGS模型时直觉选取不同视图中的同源点。实现了基于Cesium平台的Web应用，将3DGS结果转换为3D Tiles进行流式传输。支持多视图交会（超过两个视图）以提高测量精度。其设计目的不是单独增加一个模块，而是用上述表示、监督或渲染约束消除原流程在目标任务中的具体失配。提出利用3DGS渲染视图进行精确3D点位测量的新方法。利用3DGS精确渲染源视图和插值视图的能力，用户可在不同视图中直觉选取同源点，通过空间交会三角测量生成精确3D测量。实现RMSE 1-2cm精度，在薄结构和锐角区域显著优于网格测量。适用时还应保留论文给出的边界：依赖用户手动选取同源点；需要精确的相机标定；对非朗伯表面具有挑战性。

\section{其他遥感方法}

这一节的路线可以概括为“把 3DGS 变成可分割、可测量的农业场景表示”。原始 3DGS 只优化颜色和密度，最多提供一个可渲染的连续场；作物表型需要的是单个器官的长度、宽度、体积和数量。关键问题不在于再提高一点 PSNR，而在于如何把 2D 检测和分割结果稳定地提升到三维高斯实例，并从实例几何中读出生物学指标。
% METHOD: 2504.06978
\paragraph{Wheat3DGS（小麦穗实例分割与表型测量）（2025.04）}
Wheat3DGS~\cite{2504_06978} 的基线是原始 3DGS 加农业图像分割。它面对的不是普通室内物体，而是田间密集、遮挡严重、形态细长的小麦穗；MVS 点云和激光扫描要么成本高，要么在细节和速度上不适合高通量表型。它的核心改动是把 3DGS 训练出来的高斯作为三维承载体，然后把 YOLOv5 和 SAM 得到的多视图 2D 小麦穗掩码投影回高斯，求每个高斯属于哪个小麦穗实例，最后再对单个实例做形态测量。

具体流程分四步。第一步，输入未定姿的多视图田间图像，先用针对 Global Wheat Head Dataset 微调过的 YOLOv5 找出小麦穗框，再用 SAM 生成实例掩码。第二步，用 3DGS 重建作物冠层，公开代码沿用 3DGS 的训练和渲染入口，并使用支持深度、alpha 和 FlashSplat 风格贡献计算的 rasterization 组件。第三步，把每个 2D 掩码通过相机投影和渲染贡献关系映射到高斯集合，代码中的 \texttt{run\_3d\_seg.py} 体现了这一点：它累计不同视角中掩码对高斯的贡献，再用多实例优化为每个对象选择一组高斯，而不是为每个高斯训练一个额外语义网络。第四步，分割出的单穗高斯或点云进入形态测量模块，公开的 \texttt{wheatheadsmorphology} 入口会计算 PCA 主轴、样条长度、宽度、凸包体积、曲率和倾角等表型。相比 Gaussian Grouping 这类通用 3D 分割锚点，Wheat3DGS 的差异在于它不追求开放词汇编辑，而是把 2D 农业检测器、SAM 掩码、多视图高斯归属和植物形态学指标连成一条面向田间表型的管线。实验用 TLS 作几何参照，说明长度和宽度可达到表型筛选可用的量级，但体积和遮挡场景仍是更难的部分。
\textbf{实现细节：} 与锚点 TreeDGS 聚焦单棵树木的 DBH 测量不同，Wheat3DGS（2025.04） 将 3DGS 的测量对象从单一个体扩展到数百个密集排列的麦穗——问题规模从"一棵树的横截面几何"跃迁为"数百个高度相似对象的 3D 实例分割与表型提取"。这一跃迁带来核心困难：麦穗的密集排列、随机分布和外观高度相似性使得跨视角关联极其困难，现有视频追踪器（如 Gaussian Grouping 采用的 DEVA）在稀疏视角和重复结构下完全失效。为解决这一挑战，本方法提出了"检测-重建-关联-精炼-测量"的自动化流水线：（1）2D 麦穗检测与分割——使用在 Global Wheat Head Dataset 上微调的 YOLOv5 生成每张图像中每个麦穗的边界框，再逐个送入 SAM 获得精确 2D 实例掩码；（2）3DGS 场景重建——以标准 3DGS 重建整个麦田地块的 3D 场景，选择 3DGS 而非 NeRF 正是因为其更直接的场景编辑能力（后续需要对每个高斯基元赋予实例标签）；（3）跨视角实例关联——这是与 TreeDGS 的关键差异所在：TreeDGS 只需在已知位置提取横截面，不存在实例对应问题，而 Wheat3DGS 必须解决"SAM 输出的掩码是实例无关的"这一核心难题。
\textbf{实现补充：} 实现上还要核对输入、输出、参数、特征、初始化、训练损失、正则、渲染或推理接口，以及优化中的剪枝、分裂、采样、编码或解码步骤；这些抓手决定该方法到底改了 3DGS 管线的哪一环。
\textbf{方法信息：} 代码：\url{https://github.com/zdwww/wheat-3dgs}。
\section{卫星3D重建}

卫星三维重建的锚点与近景 3DGS 很不同。近景 3DGS 通常假设 pinhole 相机和同一时间段采集的照片；卫星影像则常用 RPC 相机模型，多视角之间可能跨季节、跨光照、跨时间，移动目标和辐射差异会破坏普通多视图匹配。传统 MVS 依赖监督高程或稳定纹理，NeRF 类卫星方法可以自监督但训练慢。因此该路线的核心问题是：如何把 3DGS 做成可泛化、RPC 感知、适应多时相稀疏卫星影像的前馈重建器。


% METHOD: 2508.09479
\paragraph{SkySplat（多时相稀疏卫星图像的可泛化 3DGS）（2025.08）}
SkySplat~\cite{2508_09479} 继承的是可泛化 3DGS 和卫星摄影测量两条线：一边要像 MVSPlat/HiSplat 一样从少量图像前馈预测高斯参数，另一边又不能把卫星的 RPC 几何简化成普通近景相机。它的具体改动是把 RPC 模型放进代价体和高程估计环节，用网络从多时相稀疏图像中预测高度与高斯参数，并用自监督约束绕开真实 DSM 标注。

从流程看，输入是若干幅带 RPC 的卫星图像。网络先提取多视图特征，再按 RPC 几何构造高度方向的 cost volume，估计每个像素或采样位置的高度，并据此生成可渲染的 3D Gaussians。为了处理多时相影像中的车辆、阴影、季节变化和辐射不一致，它加入 CSCM transient mask，把不稳定区域从训练重建监督中压低；同时引入预训练单目深度模型给出的相对高度作为额外约束，使高程估计不完全依赖跨时相光度一致性。最后再通过 consistency aggregation 聚合多视图预测，得到 DSM 和新视角渲染结果。项目 README 与配置说明也对应这一定位：项目不是每个场景单独优化 3DGS，而是用 PyTorch Lightning/Hydra 配置训练一个可测试的前馈模型，并明确以 RPC-aware 3DGS、自监督相对高度和快速推理为主要模块。相对 Sat-NeRF 或 EOGS 这类逐场景优化方法，SkySplat 的关键变化在于把卫星几何先验前置到网络结构中，换取跨区域推理速度；论文报告最高 86$\times$ 的速度提升，但也承认低纹理、反光区域仍继承 MVS 高程估计的弱点。
\textbf{实现细节：} SkySplat（2025.08） 针对逐场景优化卫星 3DGS 训练耗时长且无法泛化到新场景的局限，用可微 RPC 反投影代价体替代标准针孔相机假设（对每个像素在预设高程范围内采样多个候选高度，用 RPC 反投影到 3D 空间后再投影回源视角比较特征，多视角特征方差越小则该高度候选越接近真实值）、用相对高度 Pearson 监督和瞬态掩码替代纯 RGB 光度损失，核心思想是前馈式一次性预测高斯基元，比 EOGS 快 86$\times$ 且精度更高，将 DFC19（2019 IEEE GRSS Data Fusion Contest，遥感数据融合竞赛）数据集 MAE 从 13.18 m 降至 1.80 m。
\textbf{技术实现补充：} SkySplat（2025.08） 针对逐场景优化卫星 3DGS 训练耗时长且无法泛化到新场景的局限，用可微 RPC 反投影代价体替代标准针孔相机假设（对每个像素在预设高程范围内采样多个候选高度，用 RPC 反投影到 3D 空间后再投影回源视角比较特征，多视角特征方差越小则该高度候选越接近真实值）、用相对高度 Pearson 监督和瞬态掩码替代纯 RGB 光度损失，核心思想是前馈式一次性预测高斯基元，比 EOGS 快 86$\times$ 且精度更高，将 DFC19（2019 IEEE GRSS Data Fusion Contest，遥感数据融合竞赛）数据集 MAE 从 13.18 m 降至 1.80 m。
\textbf{方法信息：} 代码：\url{https://github.com/nancheng2001/skysplat-main}。
\begin{takeaway}
本章的共同脉络是：遥感和航拍中的 3DGS 很少单独作为“更漂亮的渲染器”出现，它必须与任务几何绑定。A-TDOM 把 3DGS 放进在线航测和正射图生成循环；点量测工作把 3DGS 作为高保真观察界面，再用经典多射线交会求坐标；Wheat3DGS 把 2D 检测与 SAM 掩码提升到高斯实例以服务表型测量；SkySplat 则把 RPC、多时相掩码和相对高度约束嵌入前馈 3DGS。读这些方法时，最重要的不是它们都使用 Gaussian Splatting，而是看它们把哪个遥感/测绘环节接到了高斯表示之前、之中或之后。
\end{takeaway}
\textbf{设计动机与边界：} 相对原始 3DGS 以独立高斯和图像重建误差为主的设定，SkySplat 把问题落在以下可实现的改动上：SkySplat框架包含四个主要步骤：(1) 多视角特征提取：使用多视角Transformer聚合跨视角特征，避免3D卷积；(2) RPC引导代价体构建：利用可微分RPC变换在高度假设平面上进行特征关联，通过逆RPC投影将像素坐标反投影到3D空间再投影到源视角，构建代价体；(3) 高度估计：通过soft argmin操作估计高度图；(4) 高斯参数预测：使用2D U-Net预测剩余高斯参数。CSCM模块通过一致性掩码减轻瞬态物体（移动车辆、植被变化）干扰。仅使用单目相对高度（来自预训练深度模型）作为额外监督，处理多时相图像的辐射不一致性。一致性聚合策略精化场景重建结果。其设计目的不是单独增加一个模块，而是用上述表示、监督或渲染约束消除原流程在目标任务中的具体失配。提出SkySplat自监督框架，首次将卫星RPC模型集成到可泛化3DGS流程中。包含跨自一致性模块(CSCM)处理瞬态物体干扰，以及多视角一致性聚合策略。无需真值高度图，仅用RGB图像和辐射鲁棒的相对高度监督。比EOGS快86倍且精度更高，MAE从13.18m降至1.80m。适用时还应保留论文给出的边界：RPC模型假设卫星轨道参数准确；对极端多云或雾霾场景效果可能下降；高度估计受地形起伏范围约束。
